\chapter{Analysis and Enhancement of RED2}

This chapter analyzes the performance of RED2 in the context of another
\lambdacalculus{} system and presents several ways the RED2 architecture might
be enhanced.

\section{Comparing RED2 with the Three Instruction Machine}

Insight into the efficiency of RED2 as an implementation vehicle can be gained
by comparing it to other architectures which reduce the \lambdacalculus{}.
Unfortunately, for reasons presented in Chapter 2, there is a lack of
high-performance implementations of the complete \lambdacalculus{} to use as a
basis of comparison. There is, however, a growing number of high-performance
architectures for functional languages which implement a significant portion of
the \lambdacalculus{}. Functional languages are generally used in lieu of the
\lambdacalculus{} because of their greater efficiency, so comparing RED2 to one
of these architectures would indicate if RED2 makes the complete
\lambdacalculus{} feasible for routine programming.

The functional architecture chosen for comparison is the Three Instruction
Machine~[19][20]. The Three Instruction Machine was chosen for two reasons: it
is currently considered one of the most efficient implementation techniques for
functional languages, and I had access to a Three Instruction Machine abstract
machine implementation, called WINTER~[29], and its implementors (who assisted
in my analysis). My observations are very general in nature, so a detailed
description of how the Three Instruction Machine or WINTER works is not
necessary. The reader interested in such information is referred to the
citations given above.

Table 5.1 presents measurements of the memory usage and execution times for two
benchmark programs. The sine program is a constructive mathematics approach to
programming the trigonometric sine function. It consists mainly of
\(\beta\)-substitutions and higher-order function manipulations. The game program
estimates how good a position a game player is in by looking ahead several
moves in a simple two player game. This program consists mainly of data
structure manipulations. The THOR code for both programs is listed in Appendix
A.

The benchmarks were executed on abstract machine interpreters for RED2 and
WINTER. The source code of the benchmark programs was compiled to the
instruction set of the target architecture (either RED2 or WINTER) and the
resulting machine language program was then executed by an interpreter which
simulates the target machine. Both interpreters provide monitoring facilities
for gathering run-time performance statistics. The memory usage and instruction
counts are for the simulated architecture, but the execution times indicate how
long each interpreter took to execute a benchmark.

\begin{table}[htbp]
\centering
\begin{tabular}{llrrr}
\hline
Program & Machine & Memory Usage & Instructions & Execution Time\\
\hline
sine & RED2 & 322 & 3,226 & 0.54\\
sine & WINTER & 24,887 & 10,226 & 1.87\\
game & RED2 & 1,952 & 1,274,107 & 100.0\\
game & WINTER & 58,918 & 157,224 & 12.10\\
\hline
\end{tabular}
\caption{Benchmark comparison of RED2 with WINTER, an implementation of the
Three Instruction Machine.}
\end{table}

All of the benchmark programs were run on the same Intel 80386-based personal
computer. Both interpreters were written in Microsoft C using similar coding
techniques and eight byte long simulated instruction words. Neither interpreter
has been optimized to a significant extent; the authors of both systems believe
optimizations could improve the execution speed of the interpreters a similar
amount (about 50\%).

It is difficult (and dangerous) to make comparisons based solely on the number
of instructions executed by each architecture without having actual hardware
implementations to examine. Each instruction set can be implemented in many
different ways, with a wide variation in performance. Also, the grain size of
the two instruction sets might be different. Looking only at the instruction
counts of the sine program, one might conclude that perhaps each RED2
instruction did three times the work and took three times as long to execute as
each of the WINTER instructions. The abstract interpreter execution times
provide a way to calibrate the instruction set grain sizes; the ratios of the
execution times and the ratios of instruction counts match fairly closely for
both benchmarks, indicating that the instruction sets are of comparable grain
size.

The instruction counts and execution times taken together form a legitimate
basis for comparing the relative execution speeds of the two architectures.

Two programs may seem like an absurdly small sample set upon which to base a
comparison, and in general I would agree. However, I am not making an in-depth
analysis; I am only looking for some relative indication of merit. These two
particular programs were chosen because they vividly illustrate what is good and
what is bad about the RED2 architecture.

The sine program illustrates where RED2 excels---performing the basic
operations of the \lambdacalculus{}. This comes as no surprise, because this is
exactly what RED2 was designed to do. What is surprising is the dramatic
difference in memory space required by each architecture: RED2 uses less than
one seventy-fifth of the memory used by WINTER. The primary reason for this is
that RED2 continually recycles spent memory locations, eliminating the need for
an explicit garbage collection mechanism (more will be said about this in the
next section). The version of WINTER used in these tests uses most memory
locations only once, and then must wait until a garbage collector reclaims them.
More sophisticated implementations of the Three Instruction Machine, such as
the one described in~[6], reduce the memory space consumption by more than
50\%.

The game program illustrates where RED2 performs worst---data structure
manipulation. This is unfortunate, because data structure manipulation is a
dominant feature of most functional programs. An investigation into the behavior
of the game program indicates that most of the inefficiency is caused by RED2's
current inability to share computed data structures. (The reasons for this are
presented in the next section, along with possible remedies.) This conclusion is
supported by analytical and empirical arguments.

First, the analytical argument. The game program evaluates how good a particular
move in a board game is by exploring possible future outcomes of the game if the
move is taken. This is done by generating a tree of game-boards containing all
possible future moves and then evaluating each game-board in the tree for
favorable outcomes. The tree of game-boards is created by the function
\THOR{gametree}. The \THOR{gametree} function is passed as an argument to
\THOR{evaluate}, which examines each game-board in the tree. The tree is a
computed data structure and is not shared, so each time \THOR{evaluate} asks
for a new game-board to examine, a portion of the tree must be re-created by
reducing the \THOR{gametree} function. As a result, RED2 spends most of its
time re-creating the tree of game-boards. WINTER, on the other hand, does share
computed data structures; the tree only has to be created once and can be used
as many times as needed.

\begin{table}[htbp]
\centering
\begin{tabular}{lrrr}
\hline
Program Phase & Memory Usage & Instructions & Execution Time\\
\hline
build tree & 670 & 40,210 & 2.43\\
evaluate tree & 1,176 & 93,706 & 5.80\\
total cost & 3,846 & 133,916 & 8.23\\
\hline
\end{tabular}
\caption{Experimental evidence suggesting that RED2 performed poorly on the
game program because of its current inability to share computed data
structures.}
\end{table}

Table 5.2 presents empirical data to support my conclusion. The game program
was modified so that an actual tree of game-boards was provided to
\THOR{evaluate} instead of the \THOR{gametree} function. Because the tree is no
longer a computed data structure, it can be shared. The run-time measurements
in Table 5.2 indicate that if RED2 could share computed data structures, its
performance would increase by an order of magnitude for this example.

\section{Sharing}

The performance analysis in the previous section indicates RED2 does not
perform well for some programs because it does not allow certain types of
expressions to be shared. Why is this so, and what can be done to remedy the
situation? To answer these questions requires a discussion about sharing and how
RED2 manages its memory.

When a \(\beta\)-redex is contracted, the redex's argument is substituted for
each occurrence of the redex's bound variable. One way to do this is to
substitute a separate copy of the argument for each bound variable occurrence.
This approach produces the correct result expression, but is wasteful because
each of the argument copies are reduced separately. A more efficient approach is
to reduce a single copy of the argument and share the reduced value among all
the bound variable occurrences.

Sharing of argument values occurs naturally when applicative order reduction is
used because arguments are reduced before substitution. Normal order reduction,
however, must be modified to incorporate sharing. In environment-based
implementations, the required changes are minor. The first time an argument is
accessed from the environment, it is reduced in isolation (i.e., not as part of
an application). The value of the reduced argument is substituted for the bound
variable and, if the reduced argument is closed, the environment is modified so
as to return the reduced value upon subsequent
accesses.\footnote{Open expressions cannot easily be shared because their free
variables must be adjusted properly each time the expression is substituted.
Unlike the \lambdacalculus{}, most programming languages do not allow the use
of free variables, so closedness is seldom recognized as a requirement for
sharing.} The argument must be reduced in isolation because subtle binding
errors may occur in certain cases where the argument is an abstraction which is
substituted into the head of an expression. An alternative is to reduce
arguments as part of the result graph and forbidding the sharing of arguments in
a head position.

\subsection{Why RED2 Can't Always Share}

The memory usage of RED2 during reduction is regimented and very well behaved.
In the graph memory area, memory is allocated and released only at the free
space pointer (fsp). When a reduced subgraph is no longer immediately needed
(i.e., does not form part of the final result graph), it is promptly discarded
by adjusting the fsp. Figure 5.1 illustrates how the binary arithmetic
primitives adjust the fsp after firing to discard unnecessary graph nodes. A
similar situation occurs in the environment area: memory is allocated and
released only at the environment pointer (env). When a portion of the redex
store is no longer required, it is discarded and its space reclaimed by
adjusting the env.

\begin{quote}\small
Figure 5.1: All of RED2's built-in primitive operations assist in memory
management. In this example an arithmetic primitive discards graph nodes which
are no longer accessible by adjusting the fsp. The pictures show the tip of the
result graph just before and just after the addition primitive is fired.
\end{quote}

The prompt discarding of nodes which do not form part of the final result graph
is vital to preserving the contiguous nature of subgraph spines and ensuring the
correct operation of RED2. The IF conditional primitive illustrates the need for
discarding nodes. Consider the expression
\[
\THOR{(IF }cond\ e_t\ e_f\THOR{) }arg_1\ldots arg_n.
\]
The IF selects either \(e_t\) or \(e_f\) to be applied to \(n\) arguments. For
the selected expression to be properly applied to the arguments, the IF subgraph
must be cleared from the result graph and the fsp adjusted to point at
\(arg_1\); otherwise, the arguments will not be visible to the control unit.
Figure 5.2 shows the state of the result graph just before and after the IF
primitive fires.

\begin{quote}\small
Figure 5.2: The IF primitive demonstrates why subgraphs must be promptly
discarded when no longer part of the result graph. For the arguments to be
visible to the RED2 control unit, the IF subgraph must be discarded after the
primitive is fired.
\end{quote}

Implicit in the correct implementation of sharing is the requirement that the
expression being shared exist in memory until it is no longer needed. LISP
systems~[3] and most functional language implementations meet this requirement
by storing expressions in a heap which is maintained by a garbage collector;
expressions are retained in memory until the garbage collector detects they are
no longer accessible.\footnote{Compile-time analysis can often detect when an
expression will no longer be accessible and can be safely discarded~[23][30].
Such analyses are not perfect, however, so garbage collection is still
required.} Because RED2 discards expressions which are not immediately needed as
part of the result graph, the requirement that sharable expressions be
accessible in memory cannot always be satisfied. Therefore, RED2 cannot share
many types of expressions, including computed data structures.

\subsection{What RED2 Can Share}

There are two types of expressions that are guaranteed to exist in RED2's memory
and be accessible until no longer needed, thereby opening up the possibility of
sharing. The first type are the expressions which are part of the final result
graph. In general, however, it is difficult to decide at run-time if a given
expression will be part of the final result graph. Such a decision requires
global knowledge about the current machine state and in-depth knowledge of how
each THOR primitive is implemented. As of yet, I have not found a satisfactory
solution to this problem and do not share expressions in the final result graph.

The second type of expressions are arguments that reduce to atomic data values,
such as integers or symbols. All atomic data objects are closed expressions, so
one of the two criteria for sharing is automatically met. The other criterion,
that the shared expression be accessible until it is no longer needed, is
satisfied by overwriting the argument's closure in the environment with the
reduced data value. This can be done because closures occupy two words of memory
and atomic values occupy only one word. The memory management policy of RED2
guarantees that every environment entry is accessible until it is no longer
needed.

To implement the sharing of arguments that reduce to atomic data values, minor
modifications must be made to the VAR and JOIN instructions and a new
instruction, EP (Environment Pointer), must be added:

\begin{itemize}
\item If a VAR instruction which is not at the head of a spine is executed in
the forward direction and dereferences to a closure, an EP instruction is placed
into the result graph with a data field containing the closure's address in the
environment. Figure 5.3 illustrates the effect of the modified VAR in this
situation.
\item When a JOIN instruction inserts a single instruction result into its
parent spine, a check is made to see if the insertion point is an EP
instruction. If so, the result is also written to the location indicated by the
EP's data field, overwriting the closure that was there. In this way the reduced
argument's value is shared through the environment. Figure 5.4 illustrates the
effect of the modified JOIN.
\item The EP instruction retrieves the value of the environment entry it points
to. If the value is still a closure, the closure is activated; if the value is
an atomic data object, the object is copied into the result graph.
\end{itemize}

These relatively simple modifications produce a dramatic decrease in the number
of memory cycles required for the reduction of many expressions.

\subsection{Extending RED2 to Increase Sharing}

The policy of immediately discarding subgraphs which do not form part of the
final result graph prevents RED2 from sharing closed, non-atomic expressions.
Can RED2 be changed to admit more sharing? The answer is both yes and no:
complete sharing of all sharable expressions is not possible without a major
redesign of RED2, but near-maximal sharing is within the boundaries of the
current RED2 architecture.

The discarding of subgraphs is mandatory for the current design of RED2 to work
correctly; to amend this would require significant architectural changes. Such
changes are possible, but experiments performed along these lines indicate the
resulting architecture loses much of the elegance and simplicity of RED2. In
particular, memory consumption increases significantly, approaching that of
WINTER.

A less radical approach than completely overhauling RED2 is to add a separate
memory area, the structure area, for storing data structures. Currently, RED2
treats structures the same as other expressions; however, the architecture could
be changed to treat structures differently, especially in terms of memory
management. Structures would only be created and reduced in the structure area,
and would not be discarded immediately when not part of the final result graph.
A garbage collection mechanism would be introduced to manage the structure area;
the other memory areas would still be managed by the RED2 primitives. Such
changes would allow closed data structures (including computed data structures)
to be shared.\footnote{Arguments that reduce to unbound variables can also be
shared, but their DeBruijn indices must be decremented by \(\phi\) before they
are written over the argument's closure.}

\begin{quote}\small
Figure 5.3: Forward execution of a VAR instruction that is not the head of a
spine. The variable is dereferenced, and if its value is a closure, an EP
instruction pointing to the closure is added to the result graph.
\end{quote}

\begin{quote}\small
Figure 5.4: Execution of a JOIN instruction. If the destination of the JOIN is
an EP and the reduced argument is an atomic value, the atomic value overwrites
the closure pointed to by the EP. In this way, the reduced argument is shared
through the environment.
\end{quote}

\subsection{Detecting Closed Expressions}

If RED2 is modified so that non-atomic expressions can be shared, a mechanism
for determining if an expression is closed with respect to free variables must
be incorporated into the architecture.

One way to determine if an expression is closed is to detect if its variables
``reach'' beyond the bindings in the expression. The reach of an expression is
recursively defined as follows:

\begin{itemize}
\item The reach of all constants is \(-1\).
\item The reach of a variable is equal to its DeBruijn index.
\item The reach of an application is the maximum of its operator's reach and its
operand's reach.
\item The reach of an abstraction is the reach of its body minus one.
\end{itemize}

If the reach of an expression is less than zero, the expression is closed;
otherwise, it is open. With the addition of one new register to RED2, the
calculation of an expression's reach is simple to perform during graph
traversal. The main implementation cost is that the contents of this register
must be saved on the control stack before the argument of an APP instruction is
reduced.

Experiments indicate the reach calculation increases the number of memory cycles
needed to reduce an expression by about 5\%; the exact cost varies from
expression to expression and depends on the number of APP instructions executed
during reduction.

\section{Towards Implementing \(\eta\)-Reduction}

The domain of RED2 was chosen to be the \lambdabetacalculus{} instead of the
\(\lambda\beta\eta\)-calculus for efficiency reasons. At the current time, a
complete implementation of \(\eta\)-reduction is prohibitively expensive in terms
of memory cycles; however, RED2 can be modified to implement an almost-complete
version of the \(\lambda\beta\eta\)-calculus for a fraction of the cost of a
complete version.

The \(\eta\)-reduction rule states that
\[
\lambda z.(e z)\longrightarrow e
\]
if \(z\) does not occur as a free variable in \(e\). An \(\eta\)-redex can arise
in two ways: as part of an original \(\lambda\)-expression or as a result of
\(\beta\)-reduction. An example of the first type is the expression
\[
\lambda z.((\lambda y.y)\ z),
\]
which is easily seen to be an \(\eta\)-redex. The second type of redex is more
difficult to detect; the \(\eta\)-redex in the expression
\[
\lambda z.(((\lambda y.\lambda z.z)\ z)\ z)
\]
only becomes obvious after the inner \(\beta\)-redex is contracted, yielding
\[
\lambda z.((\lambda z.z)\ z).
\]
Detecting \(\eta\)-redexes of this type sometimes requires parts of an
expression to be traversed more than once. It is this re-traversal which makes
implementing \(\eta\)-reduction so expensive.

If one is willing to forgo the detection and contraction of every
\(\eta\)-redex, RED2 can be modified to implement a nearly complete version of
the \(\lambda\beta\eta\)-calculus with about a factor of two decrease in
performance. The modified RED2 presented here may not detect every
\(\eta\)-redex because expressions are only traversed once; \(\beta\)-redexes
are contracted during the walk down the spines, and detected \(\eta\)-redexes
are contracted during the walk back up the spines.

The key to efficiently implementing near complete \(\eta\)-reduction is to
detect if an abstraction's bound variable is encountered more than once during
the traversal of the abstraction's body. If the bound variable is never
encountered during traversal, there is no possibility of the abstraction being
part of an \(\eta\)-redex. If the bound variable occurs only once, immediately
after the abstraction's LAMBDA instruction, then an \(\eta\)-redex has been
found. If the bound variable occurs more than once the abstraction may be part
of an \(\eta\)-redex, but a second traversal of the body is necessary to find
out for sure. Instead of making the second traversal, RED2 forgoes any attempt
at \(\eta\)-reducing the abstraction.

Only those abstractions which are not part of a \(\beta\)-redex are of
interest,\footnote{If an \(\eta\)-redex is the operator of a \(\beta\)-redex, it
will disappear when the \(\beta\)-redex is contracted; if it is the operand of a
\(\beta\)-redex, it may reappear in one or more places, and each copy would be
handled separately in the normal course of reduction.} so only occurrences of
variables that dereference to unbound variable objects need to be counted during
traversal. Two bits of each unbound variable (UBV) entry in the environment are
designated to record the number of times the UBV is referenced. These two bits
are called the reference bits. When a UBV entry is initially placed into the
environment area by a LAMBDA instruction, the reference bits are set to zero.
The first time a UBV is accessed, the lower reference bit is set to one; each
subsequent access sets the upper reference bit to one. In this way, the
reference bits indicate if a UBV has been referenced zero, one, or more than one
times.

The LAMBDA instruction is modified to look for \(\eta\)-redexes during reverse
execution. When a LAMBDA is executed in reverse, it checks to see if the first
instruction in the abstraction body is a non-head VAR instruction with a
DeBruijn index of 0. If this is the case, the reference bits of the LAMBDA's
unbound variable entry in the environment are checked to see if the variable
occurred only once, thereby indicating the LAMBDA and VAR instructions form an
\(\eta\)-redex. To facilitate examining the reference bits of the LAMBDA's
unbound variable entry, a pointer to the entry is pushed on the control stack
during forward execution when the UBV is first placed into the environment
area. During reverse execution, this pointer is popped off the control stack by
the LAMBDA instruction.

A detected \(\eta\)-redex is contracted by overwriting its constituent LAMBDA
and VAR instructions with a special marker, such as a NO-OP (NO OPeration)
instruction. The NO-OP instructions are ignored by the RED2 control unit, making
them effectively invisible. Figure 5.5 illustrates the detection and contraction
of an \(\eta\)-redex.

\begin{quote}\small
Figure 5.5: Contraction of an \(\eta\)-redex. When a LAMBDA instruction is
executed in reverse, it checks to see if the next instruction in memory is a
non-head VAR with a DeBruijn index of 0. If this is the case, and the reference
bits of the LAMBDA's unbound variable entry in the environment are 01, these two
instructions form an \(\eta\)-redex. The redex is contracted by overwriting the
LAMBDA and VAR instructions with no-operation instructions.
\end{quote}

When an \(\eta\)-redex is contracted, the DeBruijn index of each variable in the
redex's body must be decremented by one to compensate for the lost abstraction.
This can be done by the printing mechanism after reduction is completed.

This proposed implementation of \(\eta\)-reduction is not complete because not
all \(\eta\)-redexes are detected and contracted. The redexes in question are a
subclass of those \(\eta\)-redexes which only become visible after
\(\beta\)-reduction. Only abstractions whose bound variable is encountered once
during forward traversal are considered as candidate \(\eta\)-redexes. For
instance, the \(\eta\)-redex in the expression
\[
\lambda z.(((\lambda y.\lambda z.z)\ (z\ z))\ z)
\]
would be detected because the variable \(z\) is only encountered once during
forward traversal; the occurrences in the subexpression are never encountered.
The expression
\[
\lambda z.(((\lambda y.\lambda z.z)\ z)\ z)
\]
is an example of an \(\eta\)-redex that is not detected, because the \(z\)
variable is encountered twice during forward traversal.

The modifications described in this section have not been made a permanent part
of RED2. The degradation in performance is not worth the small gains in
reduction. Although I have no evidence beyond ``programmer's intuition'', I do
not think \(\eta\)-redexes are common or very relevant to everyday programming
tasks. More importantly, I am uneasy with the concept of ``almost complete''
reduction systems. The incompleteness weakens the system considerably from a
theoretical view.

\section{Enhanced Compilation and von-Neumann Plus Architectures}

The instruction set presented in the previous chapter is designed to reduce THOR
expressions in the most general way possible, with a minimum of compilation.
Each instruction (I am referring primarily to the primitive operations) must be
capable of handling many different possible cases. For example, the addition
primitive must deal with arguments of all kinds: integers, floats, unbound
variables, irreducible applications, etc. Implementing such generality is
expensive in terms of hardware and execution time.

There are many cases where such generality is not needed. During compilation it
is possible to detect when a more restricted form of an instruction could be
used. Arithmetic expressions are a good example. When the argument(s) of an
arithmetic operator are known to be numbers, it is not necessary to prime the
firing mechanism and traverse the argument(s) before performing the operation;
the operation can be performed immediately. These new ``immediate''
instructions treat the result graph like a value stack and are similar to the
instructions found in stack-based computers.

For instance, Figure 5.6 compares the compilation of the expression
\[
\THOR{(+ (* 2 3) (* 4 5))}
\]
using general and immediate instructions. The general instructions take 43
memory cycles to reduce the expression completely, while the immediate
instructions only take 20 memory cycles.\footnote{These numbers seem large, but
remember that no data registers are being used to cache the operands. An actual
hardware implementation of RED2 would probably include such registers, thereby
cutting the number of memory cycles required to reduce most expressions.} The
more arithmetic operations an expression contains, the greater the savings
provided by immediate instructions.

\begin{verbatim}
        (General)             (Immediate)

0: oAPP 3                0: oINT 3
1: oAPP 6                1: oINT 2
2: ePRIM_2 +             2: oIMM *
3: oINT 3                3: oINT 4
4: oINT 2                4: oINT 5
5: ePRIM_2 *             5: oIMM *
6: oINT 5                6: eIMM +
7: oINT 4
8: ePRIM_2 *
\end{verbatim}

\begin{quote}\small
Figure 5.6: A comparison of code sequences for the expression \(2 \times 3 +
4 \times 5\). The sequence on the left was generated using the general
instructions and compiler introduced in the previous chapter; the sequence on
the right was generated using more specialized, immediate instructions and a
smarter compiler.
\end{quote}

The utility of immediate instructions extends far beyond arithmetic expressions.
These instructions provide an efficient way for RED2 to implement statically
typed functional languages, which do not utilize the complete
\lambdacalculus{}. Such languages exclude the use of free (unbound) variables
and perform extensive type checking at compile time; both of these
characteristics allow a considerable increase in run-time performance.

As its instruction set expands and becomes more specialized, RED2 begins to look
like a conventional von-Neumann computer. In fact, the ideal way to implement
RED2 may be to simply add the instruction set introduced in Chapter 4 to an
existing machine. Such a hybrid, or von-Neumann Plus,\footnote{This term is
attributable to Berkling.} architecture would provide an execution platform for
both imperative and declarative programs and provides exciting possibilities for
combining the two paradigms.

Experiments performed using a microprogrammable AMD-29300 Evaluation System~[39]
indicate that such a hybrid computer is quite feasible. An instruction set for
reducing the pure \lambdacalculus{} was microcoded for the 29300 using only 125
out of 16,384 available words of microprogram memory, leaving ample space for
the inclusion of a conventional instruction set.
